% D.87 integrated after complete root and independent proof audits; historical candidate retained.
\subsection{A uniform finite-level comparison with the actual operator image}

\begin{statement}{D.87}{Proved}{The exact finite-localization kernel of the actual operator image}
For every completed pair \(s=(a,b)\) of five-adic Kummer classes
over \(\mathbb Q\), let \(I_s\) and \(I_5\) be its actual
compact global and local-at-five Alexander operator images,
including the cyclotomic multiplier, and let \(D=A_s^D\).
There is a natural identification
\[
 \mathcal S_s^*\simeq
 \ker\{H^1(I_s,D)\longrightarrow H^1(I_5,D)\}.       \tag{D.87.1}
\]
It is an identification with the filtered union of the finite
localization kernels defining \(\mathcal S_s^*\), not merely
with the unrestricted kernel of restriction to infinitely many
local groups. Every given class on the right has a representative
at one finite coefficient level which is locally trivial there
at every place.

Here is its explicit coefficient form. Let
\[
 \mathbf r=(a_5,b_5),\qquad
 L=\overline{\sum_{q\ne5}\mathbb Z_5(a_q,b_q)},\qquad
 \mathbf z=\sum_{q\ne5}\frac{\log(q^4)}{4\log6}(a_q,b_q)
                                                     \quad\in L.
                                                               \tag{D.87.2}
\]
The sums have their actual uniform \(c_0\) meaning. Write
\(f_{\mathbf v}=t_1^{v_1}t_2^{v_2}-1\) and let \(I_L\)
be the closed ideal generated by these elements for
\(\mathbf v\in L\). With
\[
 B=\Lambda/I_L,\qquad R=\Lambda/(f_{\mathbf z}),
                 \qquad J=I_LR,
                                                               \tag{D.87.3}
\]
where the relation is omitted for \(\mathbf z=0\), (D.87.1)
becomes the natural exact formula
\[
 \mathcal S_s^*\simeq
       (J^D)^{I_0}/\operatorname{im}((R^D)^{I_0}).       \tag{D.87.4}
\]
The actual \(I_0\)-action here is
\(t^{\mathbf r\kappa_5(g)}\rho_0(g)\). The formula
includes zero and nonprimitive lattices of either rank.
It does not claim an isomorphism type for every resulting
module quotient, a general section-existence theorem, or a
first higher-derived obstruction.

\end{statement}

\begin{proof}
\paragraph{Exact images and the two inclusions.}
D.83's actual closed-image and integral local-gauge argument
applies with precisely (D.87.2). It gives
\[
 I_s=L\rtimes I_0,\qquad
 I_5=\mathbb Z_5\mathbf z\rtimes I_0,                 \tag{D.87.5}
\]
where the common embedded complement is the actual graph
\(g\mapsto(\mathbf r\kappa_5(g),g)\). Its conjugation on
\(L\) is multiplication by \(\chi(g)\).
In particular the common complement is contained in \(I_5\).
These are faithful actual operator groups; no formal series
or abstract image is substituted for them.

A finite everywhere-locally-trivial class descends through
\(I_s\), by the finite affine Chebotarev argument already
proved in D.78/D.83. For completeness, if its cocycle were
nonzero on the kernel of the actual operator representation,
its finite affine image would contain \((v,1)\), \(v\ne0\).
Chebotarev supplies an unramified Frobenius in that conjugacy
class. Its linear action is the identity while its cocycle
value is a nonzero invertible linear image of \(v\), so its
local unramified class cannot vanish. This is a contradiction.
The cocycle therefore descends continuously through the compact
quotient \(I_s\). Its local-at-five class is zero on \(I_5\),
since \(G_{\mathbb Q_5}\to I_5\) is onto.
Inflation in degree one is injective. Passing to the filtered
union of finite coefficients gives the natural injection from
left to right in (D.87.1). If an element becomes zero globally
in the union, it is already zero at a later finite global
level, hence zero in that finite localization kernel; this
also checks injectivity of the union into \(H^1(G_{\mathbb Q},D)\).

For the reverse inclusion, fix a class in the right-hand kernel
and a continuous actual cocycle \(z:I_s\to D\) representing
it. Its restriction to the common graph complement \(I_0\)
is a coboundary, since \(I_0\subset I_5\). We first prove
that for every \(g\in I_s\), the element \(z(g)\) is a
\(g-1\) coboundary in the \emph{full} module \(D\).

\paragraph{The full torsion-free coefficient is essential.}
We use here \(A_s=\Lambda C\), not any nonprimitive quotient
\(R\), \(B\), or \(J\). It is a torsion-free formal power-series
module over \(\mathbb Z_5\). Let
\[
                     U_g=\chi(g)^{-1}\rho_s(g)
\]
be its compact Cartier-untwisted operator. If
\(\chi(g)\in1+5\mathbb Z_5\) is nontrivial, then
\[
                         \ker(U_g-1)=0.               \tag{D.87.6}
\]
Indeed the augmentation graded piece of degree \(n\ge0\)
is free over \(\mathbb Z_5\), and the operator on it is
multiplication by \(\chi(g)^{n+1}\): the actual Alexander
multiplier has augmentation \(\chi^2\), substitution has
weight \(\chi^n\), and the Cartier untwist removes one
factor \(\chi\). A nonzero formal series has a least
nonzero augmentation degree. Its fixedness would imply
\((\chi(g)^{n+1}-1)F_n=0\) in that torsion-free piece.
The scalar is nonzero because \(1+5\mathbb Z_5\) is
torsion-free. This is impossible. No argument on the possibly
torsion graded pieces of a cyclotomic quotient is used.

If \(\chi(g)=1\), the full operator is multiplication by
an actual unit \(\mu_g\in\Lambda^\times\). If
\(\mu_g\ne1\), then \(U_g-1\) is injective because
\(\Lambda\) is a domain. If \(\mu_g=1\), faithfulness
of the actual image, including its cyclotomic coordinate,
means that \(g\) is the identity in \(I_s\); then
\(z(g)=0\).

For every injective case just discussed, ordinary Pontryagin
duality gives surjectivity of \(g-1:D\to D\). In matrices
or operator notation this is the dual of
\(U_g^{-1}-1\), whose injectivity is equivalent to that of
\(U_g-1\). Continuous duality on compact abelian groups is
exact, so no density or rationalization is substituted for
actual surjectivity on the discrete dual.

The remaining case is \(\chi(g)\not\equiv1\pmod5\).
Write \(g=(h,g_0)\) in (D.87.5). Since
\(\chi(g)-1\in\mathbb Z_5^\times\), the translation
\(h/(\chi(g)-1)\) belongs to the actual lattice \(L\).
Conjugation by it sends \(g\) into the common complement
\(I_0\). The cocycle class is zero on that complement,
so its restriction to the cyclic subgroup generated by
\(g\) is a coboundary as well. This is invariance of
restriction under actual group conjugation; it does not
assume that the cocycle is unchanged pointwise by conjugation.
Explicitly, for a conjugating element \(b\) and
\(g'=bgb^{-1}\),
\[
 z(g')=b\,z(g)+(1-g')z(b).
\]
If \(z(g')=(g'-1)F\), then
\(z(g)=(g-1)b^{-1}(F+z(b))\).
Together these cases prove cyclic triviality for every actual
\(g\in I_s\), including all torsion cyclotomic values.

\paragraph{Compactness supplies one finite bound for all elements.}
Put
\[
       W_r=A_s/\mathfrak m^r A_s,\qquad
       D_r=W_r^D,\qquad \mathfrak m=(5,u,v).
\]
These are finite \(I_s\)-stable modules, and their increasing
union is \(D\). Continuity of \(z\) on the compact group
\(I_s\) gives finite image, so all its values lie in some
\(D_{r_0}\). For \(r\ge r_0\), define
\[
 U_r=\{g\in I_s:z(g)\in(g-1)D_r\}.                    \tag{D.87.7}
\]
Each \(U_r\) is clopen: it is the finite union, over
\(d\in D_r\), of the sets where the two continuous
finite-valued functions \(z(g)\) and \((g-1)d\) agree.
The action really factors through a finite group on this
\(I_s\)-stable coefficient module. The preceding paragraphs
show that the increasing sets \(U_r\) cover \(I_s\).
Compactness therefore gives one integer \(r\) with
\[
                              U_r=I_s.                 \tag{D.87.8}
\]
This asserts a common finite coefficient bound, not merely
existence of unrelated unbounded primitives. Continuous
selection of a primitive is not needed.

\paragraph{From the uniform cyclic bound to all local groups.}
Continuous translation coinduction in (D.87.5) permits choosing
the representative \(z\) as an inflation from
\(H^1(I_0,B^D)\), so \(z\) is zero on \(L\).
This remains a continuous cocycle in finite coefficients,
and the argument (D.87.6)--(D.87.8) applies to this representative.
The coinduction proof is valid for arbitrary closed lattices:
a Smith decomposition gives a continuous transversal for
\(\mathbb Z_5^2/L\), including finite-index and rank-zero
cases, and insertion of the identity contracts the positive
translation cochain degrees, as in D.81/D.84.

By the \(c_0\) condition, outside a finite set of primes
\(E\) away from five, the pair of inertia coordinates is
so divisible that its multiplier is trivial on \(W_r\).
Thus \(D_r\) is unramified there. The cocycle is unramified
there as well: untwisted inertia is trivial and it is zero
on the pure translation group \(L\). Apply (D.87.8) to an
actual Frobenius element. Its primitive in \(D_r\) gives
local triviality on the entire unramified local quotient.

Every prime in \(E\) where enlargement is needed has nonzero
abelian inertia. D.43 and D.80's finite local Tate-duality
argument give \(H^1(G_{\mathbb Q_q},D)=0\). Each such local
class therefore has a primitive in some finite \(D_{r_q}\).
There are only finitely many such primes. At five the assumed
vanishing on \(I_5\) gives one primitive in \(D\), hence in
some finite \(D_{r_5}\), for the whole local group. Enlarge
\(r\) to dominate these finitely many levels. All the old
primitives remain valid after inclusion, including at primes
where the enlarged coefficient has extra ramification.
The real place has no positive odd-primary cohomology.

The global cocycle is consequently locally trivial everywhere
at this one finite coefficient level. It belongs to
\(\mathcal S_s^*\), proving the reverse inclusion in (D.87.1).
No filtered union was interchanged with an infinite product
of local groups.

Finally the natural coinduction identifications in (D.87.5)
give \(H^1(I_s,D)=H^1(I_0,B^D)\) and
\(H^1(I_5,D)=H^1(I_0,R^D)\), with restriction induced by
\(B^D\subset R^D\). The exact coefficient sequence
\(0\to B^D\to R^D\to J^D\to0\) proves (D.87.4).
All nonprimitive cyclotomic components are retained in that
exact quotient. The only invariant-vanishing argument used
the full torsion-free \(A_s\), as emphasized above.
\end{proof}
